\documentclass{situationist-preprint}

\title{Two Concepts of Liberty Reconciled}
\msid{POL-09-a}
\msfamily{POL-09}
\msdiscipline{Political Philosophy}
\mstype{Principal manuscript}
\mskeywords{negative freedom; positive freedom; Berlin; MacCallum; self-mastery; interference; capacity; reduction}

\begin{document}
\maketitle

\begin{abstract}
Is freedom the absence of interference by others, or the capacity to govern one's life? Berlin distinguished the two concepts and warned against the second; MacCallum replied that there is only one concept of freedom, a triadic relation between agents, constraints and ends, and that the two traditions differ only in how they fill in its terms. We argue that freedom is both negative and positive and that neither dimension reduces to the other. MacCallum's schema is correct as an account of the form of freedom judgements, but it hides a difference of kind between two families of constraint: those imposed by other agents, and those consisting in a deficit of the agent's own capacities for self-direction. We state two conditions of adequacy and show, by contrasting cases, that reducing either dimension to the other violates one of them. The result vindicates Berlin's distinction as a distinction of dimensions, not of rival concepts.
\end{abstract}

\section{The problem}

\citet{Berlin1958} distinguished two answers to the question of political freedom. Negative freedom answers the question \enquote{What is the area within which a person is or should be left to do what he is able to do, without interference by other persons?}; positive freedom answers the question \enquote{What, or who, is the source of control that can determine someone to do this rather than that?} Berlin argued that the positive concept, by identifying freedom with the rule of a person's rational self, had licensed coercion in the name of liberty.

Two lines of reply followed. \citet{MacCallum1967} argued that the distinction rests on a confusion: every judgement of freedom has the form \enquote{$x$ is (not) free from $y$ to do or become $z$}, and the two traditions differ in what they count as agents, as constraints and as ends, not in their concept. \citet{Taylor1979} argued that negative freedom, understood as mere opportunity, cannot account for the fact that some freedoms matter more than others, or for the possibility that an agent is unfree because of internal obstacles, such as compulsions, fears or false beliefs, that prevent her from doing what she most wants; freedom is an exercise concept. On the other side, theorists of pure negative liberty have argued that only constraints imposed by other agents limit freedom, and that internal obstacles, however important, are limits of ability, not of liberty \citep{Steiner1994,Carter1999}. Republicans have argued for a third concept, freedom as non-domination \citep{Pettit1997,Skinner1998}, which the companion manuscript POL-09-b examines.

We defend a pluralist answer. Freedom has a negative dimension and a positive dimension, and neither can be reduced to the other.

\section{Framework}

\begin{definition}\label{d:dims}
Let $x$ be an agent and $z$ a range of options or ways of living. The \emph{negative dimension} of $x$'s freedom with respect to $z$ is the extent to which no other agent prevents $x$ from pursuing $z$, by physical obstruction, coercive threat or the removal of means. The \emph{positive dimension} is the extent to which $x$ has the capacities by which she can direct her own pursuit of $z$: the ability to form, assess and revise ends, and to act on her assessment rather than against it.
\end{definition}

Both dimensions fit MacCallum's schema, which we accept as an account of the form of freedom judgements; Definition~\ref{d:dims} says that its constraint variable ranges over two kinds of obstacle, what other agents do and what the agent herself lacks.

\begin{definition}\label{d:reduce}
An account of freedom \emph{reduces positive to negative freedom} if it holds that the freedom of a person with respect to $z$ is fixed by the negative dimension alone. It \emph{reduces negative to positive freedom} if it holds that the freedom of a person with respect to $z$ is fixed by the positive dimension alone.
\end{definition}

We impose two conditions of adequacy.

\begin{principle}[Prevention]\label{p:prevent}
If two persons have the same capacities for self-direction with respect to $z$, and one is prevented by others from pursuing $z$ while the other is not, the first is less free with respect to $z$ than the second.
\end{principle}

\begin{principle}[Self-direction]\label{p:self}
If two persons are equally unprevented by others with respect to $z$, and one lacks the capacity to direct her pursuit of $z$ by her own assessment while the other has it, the first is less free with respect to $z$ than the second.
\end{principle}

Principle~\ref{p:prevent} is the core of the negative tradition, and no positive theorist denies it. Principle~\ref{p:self} is the core of the positive tradition, supported by Taylor's cases; negative theorists grant that the compulsive agent lacks something, and deny only that it is freedom. Section~4 defends the principle.

\section{The argument}

Consider four persons, each with respect to the same option, leaving a house in which she has lived for years.

\begin{description}
\item[Anna] is not prevented by anyone from leaving, and is able to decide whether to leave on her own assessment of her reasons.
\item[Ben] has Anna's capacities, but the door has been locked from outside by someone who intends to keep him in.
\item[Clara] is not prevented by anyone from leaving, but suffers from a severe agoraphobia: when she decides, on reflection, that she must leave, she finds herself unable to cross the threshold.
\item[Dan] is locked in, as Ben is, and suffers from Clara's agoraphobia.
\end{description}

\begin{lemma}\label{l:indep}
The two dimensions vary independently: there are persons who agree in the negative dimension and differ in the positive, and persons who agree in the positive dimension and differ in the negative.
\end{lemma}

\begin{proof}
Anna and Clara agree in the negative dimension, since neither is prevented by anyone, and differ in the positive. Anna and Ben agree in the positive dimension and differ in the negative. Ben and Dan agree in the negative dimension and differ in the positive.
\end{proof}

\begin{theorem}\label{t:irreducible}
No account of freedom that satisfies Principles~\ref{p:prevent} and~\ref{p:self} reduces positive to negative freedom or negative to positive freedom.
\end{theorem}

\begin{proof}
Suppose an account reduces positive to negative freedom. Then Anna and Clara, who agree in the negative dimension, are equally free with respect to leaving. By Principle~\ref{p:self}, Clara is less free than Anna. So the account violates Principle~\ref{p:self}. Suppose instead that an account reduces negative to positive freedom. Then Anna and Ben, who agree in the positive dimension, are equally free. By Principle~\ref{p:prevent}, Ben is less free than Anna. So the account violates Principle~\ref{p:prevent}.
\end{proof}

\begin{corollary}\label{c:both}
On any account satisfying the two principles, a person's freedom with respect to an option depends on both dimensions, and it is maximal only if she is both unprevented and self-directing.
\end{corollary}

\begin{proof}
By Theorem~\ref{t:irreducible}, it depends on each dimension. Anna is freer than Ben, Clara and Dan, by one or both of the principles applied to the pairs of Lemma~\ref{l:indep} and to Clara and Dan, who agree in the positive dimension and differ in the negative.
\end{proof}

The argument does not require that the two dimensions be measured on a common scale, or that every pair of persons be comparable. Ben and Clara are not comparable by either principle: one is free in the first dimension and not the second, the other the reverse. We think this is correct. Whether the prisoner or the agoraphobic is less free, all things considered, depends on a weighting of the dimensions.

\section{Objections}

\begin{objection}[One concept]
MacCallum showed that there is only one concept of freedom. The present account multiplies concepts without need.
\end{objection}

\begin{reply}
We accept MacCallum's schema. The thesis is that the constraint variable in the schema ranges over two kinds of obstacle, and that the freedom of a person depends on both. MacCallum's analysis left open whether the range could be restricted to one kind; Theorem~\ref{t:irreducible} shows that it cannot without violating a principle both sides accept.
\end{reply}

\begin{objection}[Ability is not liberty]
Clara lacks an ability, not a liberty. A person who cannot swim is not unfree to swim across the river; she is unable to. Principle~\ref{p:self} confuses freedom with power \citep{Steiner1994}.
\end{objection}

\begin{reply}
The non-swimmer and Clara differ. The non-swimmer lacks a skill, and the ability to act on her assessment of her reasons is not in question: if she judged that she should learn to swim, she could. Clara's deficit is in the capacity of self-direction itself. She has assessed her reasons and concluded that she must leave, and a part of her that she does not endorse prevents her from acting on her conclusion. This is the case Taylor's argument concerns \citep{Taylor1979}: she is in the power of her compulsion. Principle~\ref{p:self} is restricted to capacities for self-direction and does not count every lack of skill as unfreedom.
\end{reply}

\begin{objection}[Internal obstacles as interference]
The positive dimension reduces to the negative if internal obstacles are treated as interferences by a part of the self, or if the lack of means is treated as prevention by those who would stop one from taking them \citep{Carter1999}.
\end{objection}

\begin{reply}
The second extension is correct and is already part of Definition~\ref{d:dims}: removing or withholding means by the threat of force is a form of prevention. The first extension abandons the negative concept. Negative freedom is freedom from the interference of other agents; if a compulsion counts as an interfering agent, the line between self and other is drawn by the agent's endorsement, and that is the line drawn by the positive tradition. The reduction succeeds only by importing the positive dimension under another name.
\end{reply}

\begin{objection}[The danger]
Berlin's warning stands: once freedom includes self-mastery, the state may claim to force people to be free by making them act as their rational selves would \citep{Berlin1958}.
\end{objection}

\begin{reply}
Berlin's argument identifies a misuse of the positive dimension: the inference from what a person's true self would choose to the permissibility of coercing her. Corollary~\ref{c:both} blocks the inference. Coercion reduces the negative dimension, and a gain in the positive dimension does not make up for it by definition, since the two are not commensurable without a weighting. Whether and when coercion is justified is a question for the theory of justice, on which POL-01 gives priority to equal basic liberties.
\end{reply}

\section{Conclusion}

Freedom judgements have one form but two dimensions. Prevention by others limits freedom, as the negative tradition holds, and a deficit in the capacity to direct oneself limits it too, as the positive tradition holds. Cases that agree in one dimension and differ in the other show that neither dimension is fixed by the other, and so that freedom is both negative and positive. The argument supplies no weighting of the dimensions, and it leaves to POL-09-b whether non-domination is a third dimension or a condition of security on the negative one.

\begin{thebibliography}{9}\small
\bibitem[Berlin(1958)]{Berlin1958} Berlin, I. (1958). \emph{Two Concepts of Liberty}. Oxford: Clarendon Press. Reprinted in \emph{Four Essays on Liberty}, Oxford University Press, 1969.
\bibitem[Carter(1999)]{Carter1999} Carter, I. (1999). \emph{A Measure of Freedom}. Oxford University Press.
\bibitem[MacCallum(1967)]{MacCallum1967} MacCallum, G.~C. (1967). Negative and positive freedom. \emph{The Philosophical Review}, 76(3), 312--334.
\bibitem[Pettit(1997)]{Pettit1997} Pettit, P. (1997). \emph{Republicanism: A Theory of Freedom and Government}. Oxford University Press.
\bibitem[Skinner(1998)]{Skinner1998} Skinner, Q. (1998). \emph{Liberty before Liberalism}. Cambridge University Press.
\bibitem[Steiner(1994)]{Steiner1994} Steiner, H. (1994). \emph{An Essay on Rights}. Oxford: Blackwell.
\bibitem[Taylor(1979)]{Taylor1979} Taylor, C. (1979). What's wrong with negative liberty. In A. Ryan (ed.), \emph{The Idea of Freedom: Essays in Honour of Isaiah Berlin}. Oxford University Press.
\end{thebibliography}

\end{document}
